% !TEX root = ../../../main.tex
% Appendix entry: catalogue.minirocket -- "MiniRocket" (Appendix A).
\entryfield{Algorithm}
The block calls \code{fingerprint} in \code{Working/Catalogue/minirocket.py}, which wraps the
\code{MiniRocket} transform of the \code{aeon} library (1.4,
\code{aeon.transformations.collection.convolution_based}), itself adapted from the authors'
code \citep{dempster2021minirocket}. Steps:
\begin{enumerate}
\item \emph{Read the windows} (\code{window_vectors}). If the window set has a
  \code{shape_file} column (a \emph{Trace shape} step ran), each window's vector is read from
  that file by \code{shape_row}. Otherwise the window is sliced from the span at its
  channel-absolute start minus $\mathrm{round}(t_0 f_s)$
  (\code{span_relative_starts} in \code{Working/types/windowset.py}). A window pool without a
  shape file is refused, naming \emph{Trace shape} as the fix.
\item \emph{Gaps} (\code{repair_gaps}): a window with at most 5\,\% non-finite samples is
  filled by linear interpolation (an end gap holds the nearest value); a window with more is
  left out and counted as \code{gaps_over_5pct}.
\item \emph{Fit set} (\code{train_mask}): the rows whose \code{split} is 0 or
  \code{"train"}, else whose \code{pool_role} is \code{"train"}. With neither column, every
  window, and a warning. A split with no training window is refused.
\item \emph{Normalise} (\code{znormalise}, when \code{normalise} is on): each window minus
  its mean, over its standard deviation; a flat window becomes zeros and is counted
  (\code{n_flat}).
\item \emph{Fit and transform} (\code{fingerprint_matrix}). \code{MiniRocket(n_kernels =
  n_features, max_dilations_per_kernel, random_state = seed)} is fitted on the fit set and
  applied to every kept window, in chunks of 4{,}096. As read from aeon's source
  (\code{docs/research/analysis-blocks-primary-sources.md}, section 12): 84 kernels of
  length 9, one for each choice of 3 of the 9 positions, with weight $-1$ everywhere and
  $+2$ at the three chosen positions; up to \code{max_dilations} exponentially spaced
  dilations per kernel; for each feature a bias, a golden-ratio quantile of the convolution
  output of one random training window; and the feature value is the proportion of positive
  values (PPV), the fraction of the convolution output above that bias. The fit learns only
  the biases and never sees a label.
\item \emph{Assemble}: the features become columns \code{mr_0000}, \code{mr_0001},
  \dots{} (float32) beside the window set's existing columns, or beside only its bookkeeping
  columns when \code{keep_upstream_features} is off. Earlier \code{mr_*} columns are
  replaced, with a warning.
\end{enumerate}
aeon rounds \code{n_features} down to a multiple of 84 (84 at least):
\code{real_feature_count}; the default 10{,}000 gives 9{,}996. aeon reseeds numpy's global
random generator inside \code{fit}; the module saves and restores it under a lock, so two
runs never interleave.

\entryfield{Assumptions}
\begin{itemize}
\item \emph{Units.} Not unit-aware. With \code{normalise} on (the default) the fingerprint
  does not depend on the unit or level; with it off it does.
\item \emph{Sampling rate.} The detectors work in samples. A 9-tap kernel at dilation $d$
  spans $8d+1$ samples, so the same shape at two sampling rates gives different
  fingerprints.
\item \emph{Window length.} At least 9 samples (aeon's floor), and one length for the whole
  set.
\item \emph{Missing values.} As in step 2: at most 5\,\% filled, otherwise left out.
\end{itemize}

\entryfield{Outputs}
A \code{WindowSet}: the kept windows (starts, length, rate unchanged) with the feature table
described above. \code{meta} holds \code{n_features} (the real count),
\code{n_features_asked}, \code{n_windows_in}, \code{n_kept}, \code{n_fit_windows},
\code{fit_on} (\code{train} or \code{all}), \code{split_column}, \code{n_left_out},
\code{n_interpolated}, \code{n_flat}, \code{seed}, \code{source} (\code{x} or
\code{shape_file}), \code{table_mb}, \code{upstream_dropped}, timings, \code{warnings},
\code{rules} and a one-line \code{summary}.

\entryfield{Cost}
$O(N \cdot L \cdot F)$ for $N$ windows of $L$ samples and $F$ features; the table is
$N \times F \times 4$ bytes. \code{max_cells} (50 million cells, 200\,MB) refuses a larger
table before any work, suggesting an \code{n_features} that fits. In a fresh process the
aeon import takes 14--45\,s and numba loads its kernels; the calibration pays these untimed.
\code{estimate} sees only the span, so it assumes the densest grid the shipped templates
make (three windows over every sample, the labels' 600/200 grid) and scales by the feature
count; a set that tiles the span costs about a third of the estimate. The transform runs on
numba's own thread count, so it also runs on a node capped at one thread.

\entryfield{Failure modes}
\begin{itemize}
\item \emph{Refusals in words}: an empty set; every window left out for gaps; windows
  shorter than 9 samples; a table over \code{max_cells}; a pool without a shape file; a
  missing or inconsistent shape file; a split with no training window.
\item \emph{No split: an optimistic downstream score.} Fitting on every window lets the
  test windows shape the biases. The block warns; the effect is bounded because no label is
  used, but it is a leak.
\item \emph{The split only governs the fit here.} \blockref{catalogue.classifier} drops the
  \code{split} column and draws its own random holdout (\blockentry{catalogue.classifier}),
  so a blocked split made upstream is not the classifier's test set.
\item \emph{Fewer features, weaker answer.} The \code{n_features} help says so: a ``no'' at
  1{,}000 features is not a ``no''.
\item \emph{Level ignored.} With \code{normalise} on, two windows differing only in level or
  size get the same fingerprint, by design.
\end{itemize}

\entryfield{Validation}
\code{tests/test_minirocket.py} checks the answer as well as the plumbing, on synthetic
windows from \code{tests/_minirocket_synthetic.py} (three shape classes that differ only in
shape, on three coloured-noise backgrounds): a random forest on the default fingerprint
scores at least 0.8 on a blocked holdout for every background; shuffled labels stay within
three binomial standard deviations of chance; labels that persist in two-hour runs with no
shape difference score at chance on the blocked split (the fingerprint does not learn time);
features of training windows are unchanged when only test windows change, and the converse.
Catch22 is measured beside it on the same split and, per the test's own note, does not lose.
\code{tests/test_template_minirocket_manual_labels.py} runs the shipped template end to end;
the template note says it beat its label-permutation null on three M2\_aug recordings
(\code{docs/prompts/blocks/reports/W2-evaluation.md}).

\entryfield{References}
\citet{dempster2021minirocket}. The algorithm was checked against the installed aeon source
in \code{docs/research/analysis-blocks-primary-sources.md}.

\entryfield{Tests}
\code{tests/test_minirocket.py} (above, plus determinism per seed, the global random state,
gap handling, the refusals, the lock, a one-thread numba run).
\code{tests/test_adapter_minirocket.py} (registration, defaults, readout, cost model and
calibration). \code{tests/test_template_minirocket_manual_labels.py} (the template).
